\documentclass[11pt]{article}
\usepackage[T1]{fontenc}
\usepackage{lmodern}
\usepackage{amsmath,amssymb,amsthm}
\usepackage[margin=1in]{geometry}
\usepackage[hidelinks]{hyperref}
\newtheorem{theorem}{Theorem}
\newcommand{\vol}{\operatorname{vol}}
\hypersetup{pdftitle={A counterexample to the Macbeath point inclusion},
  pdfauthor={Kaiwen Jin},pdfsubject={A square pyramid and its maximum symmetric intersection}}
\title{A counterexample to the Macbeath point inclusion}
\author{Kaiwen Jin\thanks{\texttt{yui.ui.siki@gmail.com}}}
\date{October 14, 2026}

\begin{document}
\maketitle
\begin{abstract}
The center of a largest centrally symmetric subset of a convex body need
not satisfy the usual dimension bound on asymmetry. We compute this center
for a square pyramid in $\mathbb R^3$ and show that reflection about it,
followed by shrinking by a factor of three, takes the apex outside the
pyramid. This gives a counterexample to Ivanov's Macbeath point inclusion
conjecture.
\end{abstract}

For a compact convex body $K\subset\mathbb R^d$ with nonempty interior,
its Macbeath point is the unique point $p$ maximizing
\[
 x\longmapsto\vol_d\bigl(K\cap(2x-K)\bigr),
 \qquad 2x-K=\{2x-y:y\in K\}.
\]
The intersection is the largest centrally symmetric subset of $K$ with
center $x$: every such subset lies in both $K$ and $2x-K$.
Ivanov asked whether this prescribed center always satisfies
$K-p\subseteq-d(K-p)$ \cite[Conjecture~2.1]{Ivanov}.
The following square pyramid disproves that assertion.

\begin{theorem}\label{thm:counterexample}
Let
\begin{equation}\label{eq:body}
 K=\bigl\{(u,v,z)\in\mathbb R^3:
  0\le z\le1,\ |u|\le1-z,\ |v|\le1-z\bigr\},
 \qquad \alpha=\frac{3-\sqrt2}{7}.
\end{equation}
Then $K$ is compact and convex with nonempty interior, and its Macbeath
point is $p=(0,0,\alpha)$. For the apex $y=(0,0,1)\in K$,
\begin{equation}\label{eq:witness}
 p-\frac{y-p}{3}
 =\left(0,0,\frac{5-4\sqrt2}{21}\right)\notin K.
\end{equation}
Consequently $K-p\nsubseteq-3(K-p)$.
\end{theorem}

\begin{proof}
The defining conditions in \eqref{eq:body} are closed linear inequalities,
and $K\subset[-1,1]^2\times[0,1]$. Thus $K$ is compact and convex.
The open box
$(-1/8,1/8)^2\times(3/8,5/8)$ lies in $K$, so its interior is nonempty.

We first show that moving a center away from the axis decreases the overlap
volume. Write a center as $x=(a,b,t)$ and put
\[
 F(a,b,t)=\vol_3\bigl(K\cap(2x-K)\bigr),\qquad V(t)=F(0,0,t).
\]
The common heights of $K$ and $2x-K$ form the interval
\[
 J_t=[0,1]\cap[2t-1,2t].
\]
At a height $z\in J_t$, their sections are squares of radii
$r_1=1-z$ and $r_2=1-2t+z$, centered at $(0,0)$ and $(2a,2b)$.
In each coordinate their intersection has length at most
$2\min\{r_1,r_2\}$. The section area is therefore at most
$4\min\{r_1,r_2\}^2$, with equality when $a=b=0$.
Integrating the section areas gives
\begin{equation}\label{eq:axial}
 F(a,b,t)\le V(t).
\end{equation}

If $0<t<1$ and $(a,b)\ne(0,0)$, the inequality is strict.
Indeed, at $z=t$ the two radii both equal $r=1-t>0$.
The intersection length in a coordinate with displacement $c$ is
$\max\{2r-2|c|,0\}$. At least one of the two lengths is strictly
less than $2r$, so the intersection area is strictly less than $4r^2$.
Intersection lengths are continuous in $z$, and $t$ is in the interior
of $J_t$. The strict area deficit thus persists on an interval of positive
length, making \eqref{eq:axial} strict after integration.
For $t\notin[0,1]$ the interval $J_t$ is empty; for $t=0$ or $1$ it
has zero length. Those centers all give zero volume.

It remains to maximize $V$. The smaller radius is $r_2$ for $z\le t$
and $r_1$ for $z\ge t$. For $0\le t\le1/2$, $J_t=[0,2t]$, so
\begin{align}
 V(t)
 &=4\int_0^t(1-2t+z)^2\,dz
   +4\int_t^{2t}(1-z)^2\,dz\notag\\
 &=\frac83\bigl((1-t)^3-(1-2t)^3\bigr).
 \label{eq:first}
\end{align}
For $1/2\le t\le1$, $J_t=[2t-1,1]$, giving
\begin{align}
 V(t)
 &=4\int_{2t-1}^t(1-2t+z)^2\,dz
   +4\int_t^1(1-z)^2\,dz\notag\\
 &=\frac83(1-t)^3.
 \label{eq:second}
\end{align}
The two formulas agree at $t=1/2$.

Since $5/4<\sqrt2<3$, we have $0<\alpha<1/4$.
On $[0,1/2]$, expanding \eqref{eq:first} yields
$V(t)=8t-24t^2+(56/3)t^3$ and the exact identity
\begin{equation}\label{eq:maximum}
 V(\alpha)-V(t)
 =\frac83(t-\alpha)^2\bigl(3+2\sqrt2-7t\bigr).
\end{equation}
The last factor is at least $2\sqrt2-1/2>0$. Hence $\alpha$ is the
unique maximizer on $[0,1/2]$.
For $1/2\le t\le1$, \eqref{eq:second} gives
$V(t)\le V(1/2)<V(\alpha)$.
Outside $[0,1]$ the volume is zero, while $V(\alpha)>0$ by
\eqref{eq:first}. Thus $\alpha$ is the unique global maximizing height.
Together with the strict comparison for nonaxial centers, this proves
that $p=(0,0,\alpha)$ is the actual unique maximizer of the prescribed
volume functional.

Finally, $y=(0,0,1)$ satisfies \eqref{eq:body}.
Its required image is the point in \eqref{eq:witness}, whose height is
negative because $4\sqrt2>5$.
Thus $y-p\notin-3(K-p)$, although $y-p\in K-p$.
\end{proof}

\paragraph{AI use.}
OpenAI Codex derived the counterexample, assisted with symbolic checks,
and drafted the exposition. The argument above is analytic and does not
depend on numerical optimization. Independent mathematical review is
requested.

\begin{thebibliography}{1}
\bibitem{Ivanov}
G.~Ivanov,
\emph{Quantitative Steinitz theorem and polarity},
arXiv:2403.14761v3 [math.MG], 10 May 2025,
\href{https://arxiv.org/html/2403.14761v3}{\texttt{arxiv.org/abs/2403.14761v3}}.
\end{thebibliography}
\end{document}
